Procesně orientované a lokálně specifické úkoly kladou důraz na sledovatelný postup (inkrementální odevzdání, mezi‑verze, laboratorní zápisníky) a na vazbu úlohy na lokální data; hodnotí se postup i výsledek, čímž se snižuje motivace ke „komoditnímu“ používání generativní AI bez dokumentace.

Praktický workflow pro návrh a provoz úlohy
\begin{enumerate}
  \item Definujte milníky: např. výběr dat, analýza, první draft, revidovaná verze, závěrečná reflexe; u každého milníku požadujte konkrétní artefakty (kód, krátký log rozhodnutí, screenshoty).
  \item Požadavek na sledovatelnost: vyžadujte meziverze, poznámky o postupu a časová razítka (exporty z OneNote nebo LMS). OneNote usnadňuje sběr mezikroků a převod rukopisu na digitální artefakty \cite{kb_ai_101_tipu_pdf_04e4a115_page_8_1}.
  \item Technické sběrné místo a validace: určete platformu/formáty (LMS, OneNote, repo; notebooky, PDF), využijte integrace Copilot/Notepad pro přípravu dílčích textů a provádějte náhodné kontroly procesních artefaktů včetně krátké metareflexe \cite{kb_ai_101_tipu_pdf_04e4a115_page_14_1, kb_ai_101_tipu_pdf_04e4a115_page_8_1}.
\end{enumerate}

Konkrétní požadované artefakty (vzorně)
\begin{itemize}
  \item První draft (kód/postup) + krátký log rozhodnutí (1--3 věty).
  \item Mezikrok: screenshot běhu, chybový výstup nebo ukázka experimentu.
  \item Revidovaný výstup s komentáři, co se změnilo a proč.
  \item Závěrečná metareflexe: použité nástroje, způsob ověření, co by se dělalo jinak (šablona viz kap. Reflexivní portfolia) \cite{kb_ai_101_tipu_pdf_04e4a115_page_8_1}.
\end{itemize}

Krátká rubrika (zkr.): dokumentace postupu a mezikroků — vysoká váha; kvalita metareflexe — vysoká váha; technická kvalita a formální náležitosti. Pozn.: zahrňte lokální data tam, kde to zvýší jedinečnost řešení (snižuje možnost přímého kopírování z veřejných zdrojů).